% AY2012 力学 解答例
% 体裁・粒度は 参考/ のPDFに揃える（行間を埋め、最終解答は \ans{} で boxed）。
% 解答・数値・問題は本年度過去問から自力で導いた（東大版は体裁の手本のみ）。
\documentclass[11pt,a4paper]{article}
\usepackage{../../../00_common/preamble}

\usetikzlibrary{decorations.pathmorphing}

\title{力学 AY2012 解答例（専門基礎科目）}
\author{}
\date{}

\begin{document}
\maketitle

% ==================================================================
\section*{第1問〔I〕}

質量 $m_A$，$m_B$ の質点 A，B を自然長 $L$・ばね定数 $k$ のばねでつなぎ，
滑らかな水平面上で $x$ 軸方向のみに運動させる。$t=0$ で質点 A に $x$ 軸正向きの
初速 $v_0$ を与える。初期に系は静止しており，平衡位置は
$x_A(0)=-\tfrac{L}{2}$，$x_B(0)=+\tfrac{L}{2}$（ばねは自然長）である。

\subsection*{(1) 運動方程式}
ばねの現在の長さは $x_B-x_A$ であり，自然長からの伸びは $(x_B-x_A)-L$ である。
伸びているとき（$x_B-x_A>L$）ばねは A を $+x$ 方向へ，B を $-x$ 方向へ引く。
復元力の符号に注意して
\begin{align*}
  &\therefore\quad m_A\ddot{x}_A = +k\bigl[(x_B-x_A)-L\bigr],\\
  &\therefore\quad m_B\ddot{x}_B = -k\bigl[(x_B-x_A)-L\bigr].
\end{align*}
\[
  \ans{\;
    m_A\ddot{x}_A = k\bigl[(x_B-x_A)-L\bigr],\quad
    m_B\ddot{x}_B = -k\bigl[(x_B-x_A)-L\bigr]\;}
\]

\subsection*{(2) $(x_B-x_A)$ の時間発展方程式}
$s\equiv x_B-x_A$ とおく。(1) の両式をそれぞれ $m_A$，$m_B$ で割って差をとると
\begin{align*}
  \ddot{s}=\ddot{x}_B-\ddot{x}_A
  &= -\frac{k(s-L)}{m_B}-\frac{k(s-L)}{m_A}
   = -k\!\left(\frac{1}{m_A}+\frac{1}{m_B}\right)(s-L).
\end{align*}
換算質量 $\mu$ を $\dfrac{1}{\mu}=\dfrac{1}{m_A}+\dfrac{1}{m_B}$，
固有角振動数を $\omega\equiv\sqrt{\dfrac{k}{\mu}}
=\sqrt{k\!\left(\dfrac{1}{m_A}+\dfrac{1}{m_B}\right)}
=\sqrt{\dfrac{k(m_A+m_B)}{m_A m_B}}$ で定めると，
$r\equiv s-L$ について
\[
  \ans{\;\ddot{r}=-\omega^2 r,\qquad
  \omega=\sqrt{\dfrac{k(m_A+m_B)}{m_A m_B}}\;}
\]
すなわち相対座標は角振動数 $\omega$ の単振動に従う。

\subsection*{(3) $t$ 秒後の $(x_B-x_A)$}
初期条件は，$t=0$ でばねが自然長だから $s(0)=L$，
よって $r(0)=0$。また $\dot{s}(0)=\dot{x}_B(0)-\dot{x}_A(0)=0-v_0=-v_0$。
$r=A\sin\omega t+B\cos\omega t$ に代入すると $B=0$，$A\omega=-v_0$ より $A=-v_0/\omega$。
\begin{align*}
  r(t)=-\frac{v_0}{\omega}\sin\omega t
  \quad\therefore\quad
  x_B-x_A = L+r = L-\frac{v_0}{\omega}\sin\omega t.
\end{align*}
\[
  \ans{\;x_B-x_A = L-\dfrac{v_0}{\omega}\sin\omega t,\qquad
  \omega=\sqrt{\dfrac{k(m_A+m_B)}{m_A m_B}}\;}
\]

\subsection*{(4) $t$ 秒後の $x_A$ と $x_B$}
外力がないので全運動量は保存し，重心は等速で動く。全質量 $M=m_A+m_B$，
重心速度と初期重心位置は
\begin{align*}
  V_G=\frac{m_A v_0}{M},\qquad
  X_G(0)=\frac{m_A(-\frac{L}{2})+m_B(\frac{L}{2})}{M}
        =\frac{(m_B-m_A)L}{2M},
\end{align*}
ゆえに $X_G(t)=X_G(0)+V_G t$。重心と相対座標 $s=x_B-x_A$ から各座標を復元すると
\begin{align*}
  x_A &= X_G-\frac{m_B}{M}\,s,\qquad
  x_B  = X_G+\frac{m_A}{M}\,s.
\end{align*}
$s=L-\dfrac{v_0}{\omega}\sin\omega t$ を代入して定数項を整理すると
（$X_G(0)-\frac{m_B}{M}L=-\frac{L}{2}$，$X_G(0)+\frac{m_A}{M}L=+\frac{L}{2}$）
\begin{align*}
  x_A(t) &= -\frac{L}{2}+\frac{m_A v_0}{M}\,t
            +\frac{m_B v_0}{M\omega}\sin\omega t,\\
  x_B(t) &= +\frac{L}{2}+\frac{m_A v_0}{M}\,t
            -\frac{m_A v_0}{M\omega}\sin\omega t.
\end{align*}
\[
  \ans{\;
  x_A=-\dfrac{L}{2}+\dfrac{m_A v_0}{M}t+\dfrac{m_B v_0}{M\omega}\sin\omega t,\quad
  x_B=+\dfrac{L}{2}+\dfrac{m_A v_0}{M}t-\dfrac{m_A v_0}{M\omega}\sin\omega t\;}
\]
（$M=m_A+m_B$，$\omega=\sqrt{k(m_A+m_B)/(m_A m_B)}$。）

\medskip
検算: $t=0$ で $x_A=-\frac{L}{2}$，$x_B=+\frac{L}{2}$，
$\dot{x}_A(0)=\frac{m_A v_0}{M}+\frac{m_B v_0}{M}=v_0$，$\dot{x}_B(0)=0$ となり，
初期条件と一致。また $m_A\dot{x}_A+m_B\dot{x}_B=m_A v_0$（一定）で運動量保存も成立。✓

% ==================================================================
\clearpage
\section*{第2問〔II〕}

質量 $m$ の放射線粒子が速さ $V$ で，原点に静止する質量 $M$ の原子核に入射する。
衝突後，粒子は $x$ 軸と $\theta$ をなす向きに速さ $v$，原子核は $x$ 軸と $\phi$ を
なす向きに速さ $u$ で進む。衝突は弾性的（運動エネルギー保存）とする。

\subsection*{(1) 保存則}
運動量の $x$，$y$ 成分とエネルギーの保存より
\begin{align*}
  &\text{運動量 }x:\quad mV = m v\cos\theta + M u\cos\phi,\\
  &\text{運動量 }y:\quad 0 = m v\sin\theta - M u\sin\phi,\\
  &\text{エネルギー}:\quad \tfrac{1}{2}mV^2 = \tfrac{1}{2}m v^2 + \tfrac{1}{2}M u^2.
\end{align*}
\[
  \ans{\;
  \begin{aligned}
    &mV = m v\cos\theta + M u\cos\phi,\quad
     0 = m v\sin\theta - M u\sin\phi,\\
    &\tfrac{1}{2}mV^2 = \tfrac{1}{2}m v^2 + \tfrac{1}{2}M u^2
  \end{aligned}\;}
\]

\subsection*{(2) 正面衝突 $(\phi=0)$ のときの $u$}
$\phi=0$ では散乱も一直線上（$\theta=0$）となり，$x$ 軸上の 1 次元弾性衝突になる。
運動量とエネルギーの保存は，散乱後の粒子速度を $v'$ として
\begin{align*}
  mV = m v' + M u,\qquad
  mV^2 = m v'^2 + M u^2.
\end{align*}
第1式より $v'=V-\dfrac{M}{m}u$。第2式へ代入すると
\begin{align*}
  mV^2 &= m\!\left(V-\frac{M}{m}u\right)^{\!2}+M u^2
   = mV^2-2MVu+\frac{M^2}{m}u^2+M u^2\\
  &\therefore\quad 0=-2MVu+M u^2\!\left(\frac{M}{m}+1\right)
   = Mu\!\left[-2V+\frac{M+m}{m}u\right].
\end{align*}
$u\neq0$ より
\[
  \ans{\;u=\dfrac{2m}{M+m}\,V\;}
\]

\subsection*{(3) $M=m$ のとき：$u$ と最大エネルギー}
$M=m$ とすると (1) の三式は質量が約せて
\begin{align*}
  V = v\cos\theta + u\cos\phi,\qquad
  0 = v\sin\theta - u\sin\phi,\qquad
  V^2 = v^2 + u^2.
\end{align*}
$\theta$ を消去するため，第1式を $v\cos\theta=V-u\cos\phi$，
第2式を $v\sin\theta=u\sin\phi$ として両辺を二乗して加えると
\begin{align*}
  v^2 &= (V-u\cos\phi)^2+(u\sin\phi)^2
       = V^2-2Vu\cos\phi+u^2.
\end{align*}
これをエネルギー保存 $v^2=V^2-u^2$ と等置して $v$ を消去すると
\begin{align*}
  V^2-u^2 &= V^2-2Vu\cos\phi+u^2
  \;\therefore\; 2u^2=2Vu\cos\phi
  \;\therefore\; u=V\cos\phi\quad(u\neq0).
\end{align*}
\[
  \ans{\;u=V\cos\phi\;}
\]
中性子から陽子（質量 $m$）に与えられるエネルギーは
$E=\tfrac{1}{2}m u^2=\tfrac{1}{2}mV^2\cos^2\phi$ であり，
$\phi=0$（正面衝突）で最大となる。
\[
  \ans{\;E_{\max}=\tfrac{1}{2}mV^2\;}
\]
これは入射運動エネルギーそのもので，等質量の正面弾性衝突では中性子が静止し，
全エネルギーが陽子へ移る（だから水素を多く含む物質は中性子の減速材として有効）。

\medskip
検算: (2) で $M=m$ とすると $u=\frac{2m}{2m}V=V$，
これは (3) で $\phi=0$ とした $u=V\cos0=V$ と一致。
また (3) のエネルギー $E=\frac12 mV^2\cos^2\phi$ は $\phi=0$ で $\frac12 mV^2$，
$\phi=\frac{\pi}{2}$ で $0$ となり，物理的に妥当。✓

\end{document}
